\chapter{古典 Kantorovich 双対と \texorpdfstring{$c$}{c}-変換}
\label{ch:seminar-06-dual}

Kantorovich 問題は線形計画であり，制約付き凸最小化として
\textbf{双対問題}（制約付き凹最大化）と対をなす．
本章では非正則化問題の双対を整備し，最適輸送の台を特徴づける
\textbf{相補性条件}と，双対を一変数に縮約する\textbf{$c$-変換}を導入する．
これは第~\ref{ch:seminar-03-entropic}・\ref{ch:seminar-04-sinkhorn}章の
エントロピー正則化・Sinkhorn が「ソフト化」していた対象の
\textbf{硬い極限}にあたり，かつ Brenier 定理（次章）の出発点となる．

本章では離散の設定（$\mathbf{a} \in \R_{\geq 0}^n$, $\mathbf{b} \in \R_{\geq 0}^m$,
$\sum_i a_i = \sum_j b_j = 1$）で議論する．
コスト行列 $\mathbf{C} \in \R_{\geq 0}^{n \times m}$ は一般のものでよい．

%% ============================================================
%% 6.1 Kantorovich 双対
%% ============================================================
\section{Kantorovich 双対}
\label{sec:sem-kantorovich-dual}

\begin{definition}{双対実行可能集合}{sem-dual-feasible}
  コスト行列 $\mathbf{C}$ に対して，\textbf{Kantorovich 双対の実行可能集合}を
  \[
    \PotentialsD(\mathbf{C})
    \defeq
    \left\{
      (\mathbf{f}, \mathbf{g}) \in \R^n \times \R^m
      \;\middle|\;
      f_i + g_j \leq C_{i,j}
      \quad (\forall\, i \in \range{n},\, j \in \range{m})
    \right\}
  \]
  で定める．$(\mathbf{f}, \mathbf{g})$ を\textbf{Kantorovich ポテンシャル}とよぶ．
  以下，行列 $(\mathbf{f} \oplus \mathbf{g})_{i,j} \defeq f_i + g_j$ と書く．
\end{definition}

\begin{proposition}{弱双対性}{sem-weak-duality-classical}
  任意の $\mathbf{P} \in \CouplingsD(\mathbf{a}, \mathbf{b})$ と
  $(\mathbf{f}, \mathbf{g}) \in \PotentialsD(\mathbf{C})$ に対して
  \[
    \inner{\mathbf{f}}{\mathbf{a}} + \inner{\mathbf{g}}{\mathbf{b}}
    \leq
    \inner{\mathbf{C}}{\mathbf{P}}.
  \]
\end{proposition}

\begin{proof}
  周辺条件 $\sum_j P_{i,j} = a_i$, $\sum_i P_{i,j} = b_j$ より
  \[
    \inner{\mathbf{f}}{\mathbf{a}} + \inner{\mathbf{g}}{\mathbf{b}}
    = \sum_i f_i \sum_j P_{i,j} + \sum_j g_j \sum_i P_{i,j}
    = \sum_{i,j} (f_i + g_j) P_{i,j}
    \leq \sum_{i,j} C_{i,j} P_{i,j}
    = \inner{\mathbf{C}}{\mathbf{P}},
  \]
  最後の不等号は $f_i + g_j \leq C_{i,j}$ と $P_{i,j} \geq 0$ による．
\end{proof}

\begin{theorem}{Kantorovich 双対定理}{sem-kantorovich-duality}
\blockmeta{space=normed; route.A=prop:sem-weak-duality-classical; route.B=prop:sem-c-transform-reduction}
  非正則化 Kantorovich 問題は双対を持ち，双対ギャップは生じない：
  \[
    \MKD_{\mathbf{C}}(\mathbf{a}, \mathbf{b})
    =
    \max_{(\mathbf{f}, \mathbf{g}) \in \PotentialsD(\mathbf{C})}
    \inner{\mathbf{f}}{\mathbf{a}} + \inner{\mathbf{g}}{\mathbf{b}}.
  \]
  右辺の最大は達成される．
\end{theorem}

\begin{proof}
  主問題を Lagrange 関数で書く．周辺制約に乗数
  $\mathbf{f} \in \R^n$, $\mathbf{g} \in \R^m$ を割り当てると
  \[
    \MKD_{\mathbf{C}}(\mathbf{a}, \mathbf{b})
    = \min_{\mathbf{P} \geq \mathbf{0}}\;
      \max_{\mathbf{f}, \mathbf{g}}
      \Bigl\{
        \inner{\mathbf{C}}{\mathbf{P}}
        + \inner{\mathbf{a} - \mathbf{P}\ones_m}{\mathbf{f}}
        + \inner{\mathbf{b} - \mathbf{P}^\top\ones_n}{\mathbf{g}}
      \Bigr\}
  \]
  である（内側の $\max$ は $\mathbf{P}$ が周辺制約を破ると $+\infty$ となり，
  満たすとき $\inner{\mathbf{C}}{\mathbf{P}}$ に等しいので，
  右辺は主問題に一致する）．
  目的関数・制約はともに線形なので，有限次元線形計画の
  \textbf{強双対定理}により双対ギャップは生じず，$\min$ と $\max$ を交換できる：
  \[
    \MKD_{\mathbf{C}}(\mathbf{a}, \mathbf{b})
    = \max_{\mathbf{f}, \mathbf{g}}
      \Bigl\{
        \inner{\mathbf{a}}{\mathbf{f}} + \inner{\mathbf{b}}{\mathbf{g}}
        + \min_{\mathbf{P} \geq \mathbf{0}}
          \inner{\mathbf{C} - \mathbf{f} \oplus \mathbf{g}}{\mathbf{P}}
      \Bigr\}.
  \]
  ここで，行列 $\mathbf{Q}$ に対し
  \[
    \min_{\mathbf{P} \geq \mathbf{0}} \inner{\mathbf{Q}}{\mathbf{P}}
    =
    \begin{cases}
      0 & (\mathbf{Q} \geq \mathbf{0}),\\
      -\infty & (\text{それ以外})
    \end{cases}
  \]
  である（ある $Q_{i,j} < 0$ なら対応する $P_{i,j} \to +\infty$ で
  $-\infty$；$\mathbf{Q} \geq \mathbf{0}$ なら $\mathbf{P} = \mathbf{0}$ が最小）．
  したがって内側の最小が $-\infty$ でない条件は
  $\mathbf{C} - \mathbf{f} \oplus \mathbf{g} \geq \mathbf{0}$，
  すなわち $(\mathbf{f}, \mathbf{g}) \in \PotentialsD(\mathbf{C})$ であり，
  そのとき値は $\inner{\mathbf{a}}{\mathbf{f}} + \inner{\mathbf{b}}{\mathbf{g}}$．
  ゆえに主張の等式を得る．主問題は実行可能かつ有界な線形計画なので
  最適解を持ち，強双対性より双対の最大も達成される
  （$c$-変換による別証明は命題~\ref{prop:sem-c-transform-reduction} を見よ）．
\end{proof}

\begin{remark}{ハード制約とそのソフト化}{sem-hard-vs-soft}
  双対の制約 $f_i + g_j \leq C_{i,j}$ は\textbf{ハード制約}である．
  エントロピー正則化の双対（命題~\ref{prop:sem-entropic-dual}）
  \[
    \max_{\mathbf{f}, \mathbf{g}}
    \inner{\mathbf{f}}{\mathbf{a}} + \inner{\mathbf{g}}{\mathbf{b}}
    - \varepsilon \sum_{i,j} e^{(f_i + g_j - C_{i,j})/\varepsilon}
  \]
  は，このハード制約を指数ペナルティで\textbf{ソフト化}したものであった
  （備考~\ref{rem:sem-soft-constraint}）．
  $\varepsilon \to 0$ では，$f_i + g_j \leq C_{i,j}$ のとき
  ペナルティ項が $0$ に，破ると $-\infty$ に向かうので，
  本章の制約付き双対が回復する．
\end{remark}

\begin{remark}{双対の解釈：集荷・配達価格}{sem-dual-interpretation}
  輸送業者が，地点 $i$ での 1 単位の\textbf{集荷}に価格 $f_i$，
  地点 $j$ への\textbf{配達}に価格 $g_j$ を課すとする．
  総請求額は $\inner{\mathbf{f}}{\mathbf{a}} + \inner{\mathbf{g}}{\mathbf{b}}$．
  依頼者がこれを受け入れる条件は，どの経路 $i \to j$ も
  自前輸送コスト $C_{i,j}$ を上回らないこと，
  すなわち $f_i + g_j \leq C_{i,j}$．
  業者はこの制約下で請求額を最大化する——これが双対問題である．
\end{remark}

%% ============================================================
%% 6.2 相補性と最適輸送の台
%% ============================================================
\section{相補性と最適輸送の台}
\label{sec:sem-complementary-slackness}

\begin{proposition}{相補性条件}{sem-complementary-slackness}
  $\mathbf{P}^\star$ を主問題の最適解，$(\mathbf{f}^\star, \mathbf{g}^\star)$ を
  双対問題の最適解とする．このとき
  \[
    P^\star_{i,j} > 0
    \;\Longrightarrow\;
    f^\star_i + g^\star_j = C_{i,j}.
  \]
  すなわち最適輸送の台は，双対制約が等号で成り立つ集合に含まれる：
  \[
    \{(i,j) : P^\star_{i,j} > 0\}
    \subseteq
    \{(i,j) : f^\star_i + g^\star_j = C_{i,j}\}.
  \]
\end{proposition}

\begin{proof}
  双対定理（定理~\ref{thm:sem-kantorovich-duality}）より
  $\inner{\mathbf{C}}{\mathbf{P}^\star}
  = \inner{\mathbf{f}^\star}{\mathbf{a}} + \inner{\mathbf{g}^\star}{\mathbf{b}}
  = \sum_{i,j}(f^\star_i + g^\star_j)P^\star_{i,j}$
  （弱双対性の証明と同じ周辺の入れ替え）．したがって
  \[
    \sum_{i,j}\bigl(C_{i,j} - f^\star_i - g^\star_j\bigr)P^\star_{i,j} = 0.
  \]
  各項は，実行可能性 $C_{i,j} - f^\star_i - g^\star_j \geq 0$ と
  $P^\star_{i,j} \geq 0$ より非負であるから，総和が $0$ なら
  すべての項が $0$．ゆえに $P^\star_{i,j} > 0$ なら
  $C_{i,j} - f^\star_i - g^\star_j = 0$．
\end{proof}

\begin{remark}{台の特徴づけと Brenier への布石}{sem-support-characterization}
  相補性は，最適輸送が「等号集合」
  $\{f^\star_i + g^\star_j = C_{i,j}\}$ の上にのみ質量を置くことを意味する．
  連続版では，この条件が
  $\supp(\pi^\star) \subseteq \{(x,y) : f^\star(x) + g^\star(y) = c(x,y)\}$
  となり，$c(x,y) = \norm{x-y}^2$ のとき
  最適 $y$ が $\partial\phi(x)$ に属する——という
  Brenier 定理の台の特徴づけへつながる（次章）．
\end{remark}

%% ============================================================
%% 6.3 c-変換
%% ============================================================
\section{\texorpdfstring{$c$}{c}-変換}
\label{sec:sem-c-transform}

双対は 2 つのポテンシャル $(\mathbf{f}, \mathbf{g})$ を持つが，
一方を他方から最適に定める操作——\textbf{$c$-変換}——により
一変数の最大化に縮約できる．

\begin{definition}{$c$-変換}{sem-c-transform-def}
  $\mathbf{g} \in \R^m$ の\textbf{$c$-変換} $\mathbf{g}^c \in \R^n$ を
  \[
    (g^c)_i \defeq \min_{j \in \range{m}} \bigl(C_{i,j} - g_j\bigr)
  \]
  で定める．同様に $\mathbf{f} \in \R^n$ に対し
  $(f^{\bar c})_j \defeq \min_{i \in \range{n}}(C_{i,j} - f_i)$ と定める．
\end{definition}

$\mathbf{g}^c$ は，$\mathbf{g}$ を固定したとき
$\PotentialsD(\mathbf{C})$ 内で $\mathbf{f}$ にとれる\textbf{最大値}である：
$f_i + g_j \leq C_{i,j}$（$\forall j$）は
$f_i \leq \min_j(C_{i,j} - g_j) = (g^c)_i$ と同値だからである．

\begin{proposition}{$c$-変換による一変数化}{sem-c-transform-reduction}
  \[
    \MKD_{\mathbf{C}}(\mathbf{a}, \mathbf{b})
    = \max_{\mathbf{g} \in \R^m}
      \Bigl\{
        \inner{\mathbf{g}^c}{\mathbf{a}} + \inner{\mathbf{g}}{\mathbf{b}}
      \Bigr\}.
  \]
\end{proposition}

\begin{proof}
  まず $(\mathbf{g}^c, \mathbf{g}) \in \PotentialsD(\mathbf{C})$：
  \[
    (g^c)_i + g_j
    = \min_k (C_{i,k} - g_k) + g_j
    \leq (C_{i,j} - g_j) + g_j
    = C_{i,j}.
  \]
  次に，任意の $(\mathbf{f}, \mathbf{g}) \in \PotentialsD(\mathbf{C})$ に対し
  $f_i \leq (g^c)_i$ であり，$\mathbf{a} \geq \mathbf{0}$ なので
  $\inner{\mathbf{f}}{\mathbf{a}} \leq \inner{\mathbf{g}^c}{\mathbf{a}}$．
  したがって $\mathbf{g}$ を固定するごとに，$\mathbf{f} = \mathbf{g}^c$ が
  目的関数を最大化する実行可能なポテンシャルである．
  定理~\ref{thm:sem-kantorovich-duality} の最大化を
  $\mathbf{f}$ について先に解いた形が主張の式である．
\end{proof}

\begin{remark}{$c$-凹関数と Legendre 変換}{sem-c-concave}
  $\mathbf{g}^c$ の形をしたベクトル（アフィン関数
  $j \mapsto C_{i,j} - g_j$ の $\min$）を\textbf{$c$-凹}とよぶ．
  最適ポテンシャルは $c$-凹なものにとれるため，
  双対は $c$-凹ポテンシャルの有界集合上の最大化に帰着し，
  最大が達成される（定理~\ref{thm:sem-kantorovich-duality} の達成性）．
  $C_{i,j} = \norm{x_i - y_j}^2$ の場合，$c$-変換は
  凸関数の \textbf{Legendre 変換}と対応し，これが Brenier 定理で
  最適写像を凸関数の勾配として特徴づける鍵となる（次章）．
\end{remark}

%% ============================================================
%% 6.4 Sinkhorn との対応：ソフト $c$-変換
%% ============================================================
\section{Sinkhorn との対応：ソフト \texorpdfstring{$c$}{c}-変換}
\label{sec:sem-soft-c-transform}

エントロピー正則化（第~\ref{ch:seminar-03-entropic}章）と
Sinkhorn（第~\ref{ch:seminar-04-sinkhorn}章）は，
この $c$-変換を\textbf{ソフト化}した反復として理解できる．
$\min$ をなめらかにした\textbf{ソフト最小値}を導入する．

\begin{definition}{ソフト最小値}{sem-soft-min-ch6}
  $\mathbf{z} \in \R^m$ と $\varepsilon > 0$ に対し
  \[
    \smin_\varepsilon(\mathbf{z})
    \defeq
    -\varepsilon \log \sum_{j=1}^m e^{-z_j/\varepsilon}.
  \]
  これは $\min_j z_j - \varepsilon\log m \leq \smin_\varepsilon(\mathbf{z})
  \leq \min_j z_j$ を満たし，$\varepsilon \to 0$ で
  $\smin_\varepsilon(\mathbf{z}) \to \min_j z_j$ となる．
\end{definition}

\begin{definition}{ソフト $c$-変換}{sem-soft-c-transform}
  $\mathbf{g} \in \R^m$ の\textbf{ソフト $c$-変換}
  $\mathbf{g}^{c_\varepsilon} \in \R^n$ を
  \[
    (g^{c_\varepsilon})_i
    \defeq
    \smin_\varepsilon\bigl((C_{i,j} - g_j)_j\bigr)
    = -\varepsilon \log \sum_j e^{-(C_{i,j} - g_j)/\varepsilon}
  \]
  で定める．$\varepsilon \to 0$ のとき $\mathbf{g}^{c_\varepsilon} \to \mathbf{g}^c$．
\end{definition}

\begin{proposition}{Sinkhorn はソフト交互 $c$-変換である}{sem-sinkhorn-soft-c}
  対数領域変数 $f_i = \varepsilon\log u_i$, $g_j = \varepsilon\log v_j$
  で表すと，Sinkhorn の $\mathbf{u}$ 更新
  $\mathbf{u} = \mathbf{a} \oslash (\mathbf{K}\mathbf{v})$
  （第~\ref{ch:seminar-04-sinkhorn}章）は
  \[
    f_i \leftarrow (g^{c_\varepsilon})_i + \varepsilon\log a_i
  \]
  と書ける．$\varepsilon \to 0$ のとき，$\varepsilon\log a_i \to 0$
  （$a_i > 0$ は固定）なので，これは古典的 $c$-変換
  $f_i \leftarrow (g^c)_i = \min_j(C_{i,j} - g_j)$ に収束する．
\end{proposition}

\begin{proof}
  $u_i = a_i / (\mathbf{K}\mathbf{v})_i$ の両辺の対数を $\varepsilon$ 倍すると
  $f_i = \varepsilon\log a_i - \varepsilon\log(\mathbf{K}\mathbf{v})_i$．
  Gibbs カーネル $K_{i,j} = e^{-C_{i,j}/\varepsilon}$ と
  $v_j = e^{g_j/\varepsilon}$ より
  \[
    (\mathbf{K}\mathbf{v})_i
    = \sum_j e^{-C_{i,j}/\varepsilon} e^{g_j/\varepsilon}
    = \sum_j e^{-(C_{i,j} - g_j)/\varepsilon},
  \]
  ゆえに $-\varepsilon\log(\mathbf{K}\mathbf{v})_i = (g^{c_\varepsilon})_i$．
  したがって $f_i = (g^{c_\varepsilon})_i + \varepsilon\log a_i$．
  $\varepsilon \to 0$ で $(g^{c_\varepsilon})_i \to (g^c)_i$ かつ
  $\varepsilon\log a_i \to 0$（$a_i > 0$ 固定）なので
  $f_i \to (g^c)_i$．
\end{proof}

\begin{remark}{本章のまとめと展望}{sem-ch6-summary}
  非正則化 Kantorovich 問題は双対
  $\max_{(\mathbf{f},\mathbf{g}) \in \PotentialsD(\mathbf{C})}
  \inner{\mathbf{f}}{\mathbf{a}} + \inner{\mathbf{g}}{\mathbf{b}}$ を持ち
  （定理~\ref{thm:sem-kantorovich-duality}），
  最適輸送の台は相補性条件で特徴づけられた
  （命題~\ref{prop:sem-complementary-slackness}）．
  $c$-変換により双対は一変数化され
  （命題~\ref{prop:sem-c-transform-reduction}），
  Sinkhorn はそのソフト化（ソフト $c$-変換）の交互適用として
  古典双対に接続する
  （命題~\ref{prop:sem-sinkhorn-soft-c}）．

  次章では地の空間を $\R^d$，コストを $\norm{x-y}^2$ とし，
  $c$-変換が Legendre 変換に，相補性条件が
  「最適写像は凸関数の勾配 $T = \nabla\phi$」という
  \textbf{Brenier の定理}に結実することを見る．
  これが測地線・Wasserstein 幾何の出発点となる．
\end{remark}
